\section{The correcting twist and the quadratic field of the Case III seed}

Two arithmetic choices enter the tensor construction of an integral seed.
The scalar twist must correct the multiplier integrally, with prescribed
first and last residual characters. The quadratic field used for the
dihedral factor must remain visible after passing to the kernel of the
full cyclotomic character on every later CM tower. Both assertions can be
proved before constructing the required automorphic representations.

% BEGIN E65T1
\paragraph{E65T1. An integral square root and the labelled determinant identity.}
Fix the original Case III data and any tower $T_0$, with its Boolean label.
Write $O=c.O$, $\varepsilon$ for the integral cyclotomic character,
$\omega$ for its Teichm\"uller part, and $\psi,\xi$ for the two finite
characters of the datum. All the characters in this argument are defined
on the absolute Galois group of $\mathbb Q$. Put
\[
 \mu=\varepsilon\omega^{-1},\quad P=\omega^e\psi,\quad
 \alpha=\mu\xi,\quad \Psi=\varepsilon^e\psi,\quad
 \Theta=\alpha\Psi .
\]
We have $\overline\mu=1$, $\overline\Psi=\overline P$, and
$\ell-1\mid 2e$ in Case III (TFT1).
For each $g$, Hensel's lemma in the complete coefficient DVR gives a unique
$\eta(g)\equiv1\pmod{\mathfrak m_O}$ such that $\eta(g)^2=\mu(g)$:
the derivative of $X^2-\mu(g)$ at $1$ is a unit since $\ell$ is odd.
Uniqueness makes $\eta$ multiplicative. For any $g,h$, the sum
$\eta(g)+\eta(h)$ is a unit; factoring the difference of squares shows
that congruence of $\mu(g),\mu(h)$ modulo $\mathfrak m_O^n$ implies
the same congruence of their roots. Thus $\eta$ is adically continuous.

Let $k\ge1$ with $k-1\equiv e\pmod{\ell-1}$ and set $j=e+2-k\in\mathbb Z$.
For label false put $t=1$, $\varphi=\xi\psi$; for label true put
$t=P$, $\varphi=\xi\psi^{-1}$. In both cases put
$\Phi=\varepsilon^{k-1}\varphi$ and $\delta=t\eta^j$. Then
\[
                           \delta^2\Phi=\Theta .
\]
Indeed $\eta^{2j}=\varepsilon^j\omega^{-j}$, $j+k-1=e+1$ and
$j\equiv1\pmod{\ell-1}$. In the false case the product is
$\varepsilon^{e+1}\omega^{-1}\xi\psi$ as required. In the true case
$t^2\xi\psi^{-1}=\omega^{2e}\xi\psi=\xi\psi$ since $\omega^{\ell-1}=1$.
These formulas also establish continuity of the twist, including when $j<0$.
% END E65T1

% BEGIN E65T2
\paragraph{E65T2. The residue, kernel and weight conditions of the twist.}
Write $\chi_i$ for the four integral characters in the prescribed order.
For the false label they are $(\Theta,\Psi,\alpha,1)$ and for the
true label they are $(\alpha,1,\Theta,\Psi)$. Thus
$\chi_0\chi_3=\Theta$ in both orders and $\overline{\chi}_3=\bar t=\bar\delta$.
Reducing the identity of E65T1 and cancelling this unit gives
$\overline{\chi}_0=\bar\delta\,\bar\Phi$.

Suppose the full cyclotomic character at $g$ is $1$ and the original
residual diagonal representation equals the identity there. Then
$\omega(g)=1$ since it lifts the residual cyclotomic value and has
order dividing $\ell-1$. Reading the diagonal at the two indices
corresponding to $\alpha$ and $\Psi$ gives
$\bar\xi(g)=\bar\psi(g)=1$. Both finite characters here have
order prime to $\ell$. Reduction is injective on such torsion:
if $u\equiv1$ and $u^m=1$ for $\ell\nmid m$, the second factor of
$(u-1)(1+u+\cdots+u^{m-1})=0$ is a unit. It follows that
$\xi(g)=\psi(g)=1$. Also $\mu(g)=1$, hence $\eta(g)=1$ by the
unique square root condition. Therefore $\delta(g)=\varphi(g)=1$.

The subgroup $U=\ker\omega\cap\ker\psi$ is open. On $U$ we have
$t=1$ and $\eta^2=\varepsilon$, so $\delta^2=\varepsilon^j$ there.
This proves every clause of $\mathrm{Tw4}(T_0,k,\delta)$, for every
$T_0$ and $k$ in the source assertion $\mathrm{TwEx4}$.
% END E65T2

% BEGIN E65T3
\paragraph{E65T3. The seed's quadratic character on the full cyclotomic kernel.}
Fix seed data $d$ as in the original $\mathrm{AlNe}$ hypothesis, and
any acceptable later tower $T\ge T_0$. Write $F=T.F$, $L=T.L$, and
$K'$ for the quadratic field in $d$. On $G_F$ the residual root at
$w_\alpha$ is the character $s$ of $K'$ restricted to $G_F$: the
initial $\mathrm{QuadOn}$ asserts that it is $1$ on $G_F\cap G_{K'}$
and $-1$ elsewhere, and restriction to a later tower preserves the label
and both root ratios. Restrict $s$ further to $G_L$. This quadratic
character is continuous because its kernel is open.

Suppose it were trivial on $G'=G_L\cap\ker\varepsilon$. It would
factor continuously through $H=\varepsilon(G_L)\subseteq\mathbb Z_\ell^\times$:
the continuous surjection from the compact group $G_L$ to the Hausdorff
group $H$ is a quotient map. The closed subgroup
$H\cap(1+\ell\mathbb Z_\ell)$ is pro-$\ell$. Indeed its images modulo
$\ell^n$ are finite $\ell$-groups and these quotients separate points.
Every continuous character of this subgroup of order dividing two is
trivial since $\ell$ is odd. Consequently $s(g)=1$ for every
$g\in G_L$ with $\omega(g)=1$.

But $s\bar\chi_1=\bar\chi_0$ and $s\bar\chi_3=\bar\chi_2$,
by the definitions of $w_\alpha$ and of either ordering in E65T2.
Since $s^2=1$, multiplying the four diagonal characters by $s$ effects
the permutation $(0\,1)(2\,3)$ on $G_L$. The quadratic-self-twist
clause of $T.\mathrm{Ok}$ now forces $s|_{G_L}=1$. The root
nontriviality clause of that same record supplies $h\in G_L$ at
which the $w_\alpha$ root differs from $1$, a contradiction.
Thus some $g\in G'$ acts nontrivially on $K'$. This supplies
$\mathrm{AlNe}$ after every acceptable tower extension, even when
the marked element itself has nontrivial full cyclotomic value.
% END E65T3
